\section{狭义相对论}
\sourceproblem{选择1}{63}{69}
\problemtext 判断真空光速与频率及运动关系、各惯性系光速各向同性、惯性系物理定律等价三个命题。
\workings 三者均为狭义相对论基本原理或其直接含义，选D。
\crosscheck 仅限真空光速和惯性参考系；介质中的光传播速度另受介质影响。
\sourceproblem{选择2}{63}{69}
\problemtext 比较同地同时事件及异地同时事件在其他惯性系中的同时性。
\workings 洛伦兹变换 $\Delta t'=\gamma(\Delta t-\bm v\cdot\Delta\bm r/c^2)$。同地同时给 $\Delta t'=0$；异地同时一般不同时，教材选A。
\crosscheck 若相对速度恰垂直于两事件空间位移，仍可同时；故第二项严格表述为不一定同时。
\sourceproblem{选择3}{63}{69}
\problemtext K系沿x相隔3 m两事件同时，K'系空间间隔5 m，求相对速率。
\workings $\Delta x'=\gamma\Delta x$，故 $\gamma=5/3$，$v=c\sqrt{1-1/\gamma^2}=4c/5$，选A。
\crosscheck 此处比较事件空间坐标差，不是同一杆的长度测量。
\sourceproblem{选择4}{63}{69}
\problemtext K'系同地事件固有时间间隔 $\Delta t_0$，K系测得 $\Delta t$；K'静止杆长 $l_0$ 沿运动方向，K系长度 $l$，比较大小。
\workings 对非零相对速度，$\Delta t=\gamma\Delta t_0>\Delta t_0,\ l=l_0/\gamma<l_0$，选D。
\crosscheck K系事件发生于不同位置，须用同步钟网络测坐标时间；杆长度须在K系同时测端点。
\sourceproblem{选择5}{63}{69}
\problemtext 边长a的静止正方形位于xy面，观察系沿x以0.6c运动，求测得面积。
\workings 沿x边长收缩为 $a\sqrt{1-0.6^2}=0.8a$，y边不变，面积 $0.8a^2$，选B。
\crosscheck 这是同时坐标测量的面积，而非考虑光传播时延的照片外观。
\sourceproblem{选择6}{63}{69}
\problemtext 固有长100 m的飞船以0.6c飞行，光脉冲由尾至头，求地面系发射和接收事件空间间隔（选项续PDF70）。
\workings 船系 $\Delta x'=100$ m、$\Delta t'=100/c$。逆变换 $\Delta x=\gamma(\Delta x'+v\Delta t')=1.25(100)(1.6)=200$ m，选C。
\crosscheck 地面系船长80 m，追及时间 $80/(c-0.6c)=200/c$，与变换一致。
\sourceproblem{选择7}{64}{70}
\problemtext 年发电量 $10^{11}$ kWh，若全部来自材料静能转化，求所需质量。
\workings $E=10^{11}(3.6\times10^6)=3.6\times10^{17}$ J，$m=E/c^2\approx4.0$ kg，选A。
\crosscheck 这是百分之百静能转化的假设值，不是实际核燃料消耗量。
\sourceproblem{选择8}{64}{70}
\problemtext 电子动能0.25 MeV、静能0.51 MeV，求相对论速率。
\workings $\gamma=1+0.25/0.51=76/51$，$v/c=\sqrt{1-(51/76)^2}\approx0.741$，最接近0.75，选C。
\crosscheck 不采用非相对论动能公式；速率小于c。
\sourceproblem{填空1}{64}{70}
\problemtext 表述狭义相对论的两条基本原理。
\workings 相对性原理：物理基本定律在所有惯性系中形式相同。光速不变原理：真空中光在所有惯性系各方向的传播速率均为c，与光源及观察者的运动无关。
\crosscheck 不扩展到任意非惯性坐标系的全局坐标光速。
\sourceproblem{填空2}{64}{70}
\problemtext 相对地球匀速运动的恒星发出的光子，相对地球的速率。
\workings 真空光速恒为 $c$；沿运动方向可由速度变换 $(c+v)/(1+vc/c^2)=c$ 验证。
\crosscheck 光的频率可因多普勒效应而改变，真空速率不变。
\sourceproblem{填空3}{64}{70}
\problemtext 船头发光至船尾，船系钟测得传播时间 $\Delta t$，求固有长度。
\workings 船系端点静止，光速c，故 $l_0=c\Delta t$。
\crosscheck 此时间为船系同步钟给出的坐标间隔，不是两事件的固有时（光状间隔固有时为零）。
\sourceproblem{填空4}{64}{70}
\problemtext 运动米尺沿长度方向测得0.8 m，求速率。
\workings $0.8=\sqrt{1-v^2/c^2}$，故 $v=0.6c\approx1.8\times10^8$ m/s。
\crosscheck 米尺固有长1 m，测端点同时性以观察系为准。
\sourceproblem{填空5}{64}{70}
\problemtext 固有长 $l_0>l$ 的杆沿长度方向运动，在门系同时容入宽l的门，求最低速率。
\workings 由 $l_0\sqrt{1-v^2/c^2}\le l$，得到 $v_{\min}=c\sqrt{1-(l/l_0)^2}$。
\crosscheck 这是门系同时测量的几何条件，不表示杆系两端同时通过，也不涉及碰撞动力学。
\sourceproblem{填空6}{64}{70}
\problemtext 地球系路程5光年，宇航员希望船系路程为3光年，求匀速段速率。
\workings $3/5=\sqrt{1-v^2/c^2}$，故 $v=0.8c$。
\crosscheck 地球系历时6.25年，船系固有时3.75年；忽略起停加速段。
\sourceproblem{填空7}{64}{70}
\problemtext 按原教材相对论质量记号，写质量与速度及动能关系。
\workings 若定义相对论质量 $m=\gamma m_0$，则 $m=m_0/\sqrt{1-v^2/c^2}$，$K=(\gamma-1)m_0c^2$。
\crosscheck 现代惯例通常使用不变质量 $m_0$，速度依赖的是能量和动量，而非不变质量。
\sourceproblem{填空8}{64}{70}
\problemtext 静质量 $m_e$ 的电子从静止加速到0.6c，求外界净功。
\workings $\gamma=1.25$，故 $W=\Delta K=0.25m_ec^2$。
\crosscheck 若有辐射或其他能量损失，施力装置供能需另计；所求为粒子净功。
